\chapter{Pecahan: Membandingkan dan Menjumlahkan}\label{ch:g7:fractions}

Kelas 6 memperkenalkan \omterm{def:g6:fractions:def}{pecahan} sebagai bagian dan sebagai \omterm{def:g3:division:remainder}{hasil bagi}
yang tepat (\cref{ch:g6:fractions}). Bab ini belajar \emph{menghitung}
dengannya: mengenali \omterm{def:g6:fractions:def}{pecahan} yang senilai, membandingkan, menjumlahkan
dan mengurangkan --- untuk sementara dengan \omterm{def:g4:fractions:def}{penyebut} yang ramah; hal
umumnya menyusul pada \cref{ch:g9:fractions}.

\section{Pecahan yang senilai}

\begin{theorem}[Pecahan senilai]\label{thm:g7:fractions:equal}
Sebuah \omterm{def:g6:fractions:def}{pecahan} tidak berubah ketika \omterm{def:g4:fractions:def}{pembilang} dan \omterm{def:g4:fractions:def}{penyebutnya}
dikalikan, atau dibagi, dengan bilangan bukan \omterm{def:g1:counting:zero}{nol} yang sama:
\[
\frac{a}{b} = \frac{a \times k}{b \times k} .
\]
\end{theorem}
\admitted

\begin{example}\label{ex:g7:fractions:equal}
$\dfrac{3}{4} = \dfrac{3 \times 5}{4 \times 5} = \dfrac{15}{20}$, dan
ke arah sebaliknya $\dfrac{42}{30} = \dfrac{42 \div 6}{30 \div 6}
= \dfrac{7}{5}$ (disederhanakan dengan $6$).

Untuk menguji apakah $\dfrac{8}{12}$ dan $\dfrac{10}{15}$ senilai,
sederhanakan keduanya: $\dfrac{8}{12} = \dfrac{2}{3}$ dan
$\dfrac{10}{15} = \dfrac{2}{3}$ --- ya, senilai.
\end{example}

\section{Membandingkan pecahan}

\begin{method}[Membandingkan dua pecahan]\label{met:g7:fractions:compare}
\begin{enumerate}
  \item \emph{\omterm{def:g4:fractions:def}{Penyebutnya} sama}: bandingkan \omterm{def:g4:fractions:def}{pembilangnya}:
        $\frac{5}{9} < \frac{7}{9}$.
  \item \emph{Salah satu \omterm{def:g4:fractions:def}{penyebut} adalah \omterm{def:g5:division:multiple}{kelipatan} \omterm{def:g4:fractions:def}{penyebut} yang
        lain}: tulislah ulang \omterm{def:g6:fractions:def}{pecahan} yang lebih kasar, lalu
        bandingkan: $\frac{5}{6}$ \omterm{def:g7:negatives:def}{lawan} $\frac{7}{12}$:
        $\frac56 = \frac{10}{12} > \frac{7}{12}$.
  \item \emph{Bandingkan dengan $1$ atau $\frac12$} bila mungkin:
        $\frac{9}{8} > 1 > \frac{7}{9}$ langsung menyelesaikan
        $\frac98$ \omterm{def:g7:negatives:def}{lawan} $\frac79$.
\end{enumerate}
\end{method}

\begin{omfigure}
\begin{tikzpicture}[scale=0.62]
  \foreach \i in {0,...,5} \draw[omExo] (\i*2,0) rectangle (\i*2+2,1.1);
  \foreach \i in {0,...,4} \fill[omDef!35] (\i*2+0.05,0.05)
    rectangle (\i*2+1.95,1.05);
  \draw[thick] (0,0) rectangle (12,1.1);
  \node[right, font=\small] at (12.3,0.55) {$\frac56$};
  \begin{scope}[yshift=-1.9cm]
  \foreach \i in {0,...,11} \draw[omExo] (\i,0) rectangle (\i+1,1.1);
  \foreach \i in {0,...,6} \fill[omThm!35] (\i+0.05,0.05)
    rectangle (\i+0.95,1.05);
  \draw[thick] (0,0) rectangle (12,1.1);
  \node[right, font=\small] at (12.3,0.55) {$\frac{7}{12}$};
  \end{scope}
\end{tikzpicture}

{\small Membandingkan $\frac56$ dan $\frac{7}{12}$ pada dua batang
kembar: memotong perenaman menjadi dua memperlihatkan
$\frac56 = \frac{10}{12}$, lebih besar daripada $\frac{7}{12}$.}
\end{omfigure}

\section{Menjumlahkan dan mengurangkan}

\begin{theorem}[Penyebut yang sama]\label{thm:g7:fractions:add}
\omterm{def:g6:fractions:def}{Pecahan} dengan \omterm{def:g4:fractions:def}{penyebut} yang sama dijumlahkan (atau dikurangkan)
dengan menjumlahkan (atau mengurangkan) \omterm{def:g4:fractions:def}{pembilangnya}:
\[
\frac{a}{d} + \frac{b}{d} = \frac{a + b}{d},
\qquad
\frac{a}{d} - \frac{b}{d} = \frac{a - b}{d} .
\]
\end{theorem}

\begin{proof}[Mengapa]
$a$ bagian berukuran $\frac1d$ ditambah $b$ bagian dengan ukuran yang
sama menjadi $a + b$ bagian berukuran itu.
\end{proof}

\begin{method}[Penyebut yang berbeda]\label{met:g7:fractions:add}
Ketika satu \omterm{def:g4:fractions:def}{penyebut} merupakan \omterm{def:g5:division:multiple}{kelipatan} \omterm{def:g4:fractions:def}{penyebut} yang lain:
\begin{enumerate}
  \item tulislah ulang \omterm{def:g6:fractions:def}{pecahan} yang \omterm{def:g4:fractions:def}{penyebutnya} lebih kecil agar
        keduanya berpenyebut yang lebih besar itu
        (\cref{thm:g7:fractions:equal});
  \item jumlahkan atau kurangkan \omterm{def:g4:fractions:def}{pembilangnya};
  \item sederhanakan hasilnya kalau bisa.
\end{enumerate}
\end{method}

\begin{example}\label{ex:g7:fractions:add}
Hitunglah $\dfrac{3}{4} + \dfrac{5}{8}$, selangkah demi selangkah:
\begin{enumerate}
  \item $8$ adalah \omterm{def:g5:division:multiple}{kelipatan} $4$: tulis ulang
        $\dfrac34 = \dfrac{3 \times 2}{4 \times 2} = \dfrac{6}{8}$;
  \item jumlahkan: $\dfrac{6}{8} + \dfrac{5}{8} = \dfrac{11}{8}$;
  \item tidak ada yang perlu disederhanakan: hasilnya
        $\dfrac{11}{8}$, yaitu $1 + \dfrac38$.
\end{enumerate}
Satu lagi: $2 - \dfrac{4}{5} = \dfrac{10}{5} - \dfrac{4}{5} =
\dfrac{6}{5}$ (tulislah dulu bilangan bulatnya di atas \omterm{def:g4:fractions:def}{penyebut}
$5$).
\end{example}

\begin{example}[Perangkap yang klasik]\label{ex:g7:fractions:trap}
$\dfrac12 + \dfrac13$ itu \emph{bukan} $\dfrac{2}{5}$! Menjumlahkan
\omterm{def:g4:fractions:def}{pembilang} dan \omterm{def:g4:fractions:def}{penyebut} secara terpisah itu salah --- jawaban
$\frac25$ akan lebih kecil daripada $\frac12$, dan itu mustahil
ketika kita \emph{menjumlahkan} sesuatu yang positif pada $\frac12$.
(\omterm{def:g1:addition:def}{Jumlah} yang benar, $\frac56$, memerlukan \omterm{def:g4:fractions:def}{penyebut} bersama $6$; cara
umumnya ada pada \cref{ch:g9:fractions}.)
\end{example}

\section{Mengalikan pecahan dengan sebuah bilangan}

\begin{proposition}[Pecahan dikali bilangan]\label{prop:g7:fractions:mult}
Untuk sebuah bilangan $k$:
\[
k \times \frac{a}{b} = \frac{k \times a}{b} .
\]
Mengambil $\frac ab$ \emph{dari} suatu besaran berarti mengalikannya
dengan $\frac ab$.
\end{proposition}
\admitted

\begin{example}\label{ex:g7:fractions:mult}
$6 \times \dfrac{2}{3} = \dfrac{12}{3} = 4$: enam kali dua pertiga
adalah empat utuh. Dan ``$\frac34$ dari $60$ euro'' adalah
$\frac34 \times 60 = \frac{180}{4} = 45$ euro --- jawaban yang sama
seperti membagi dengan $4$ lalu mengalikan dengan $3$
(\cref{met:g6:fractions:ofquantity}).
\end{example}

\section{Latihan}

\begin{exercise}[$\star$]\label{exo:g7:fractions:1}
Lengkapilah: $\dfrac{2}{5} = \dfrac{?}{20}$; \quad
$\dfrac{18}{24} = \dfrac{3}{?}$; \quad
$\dfrac{7}{3} = \dfrac{28}{?}$; \quad
$\dfrac{?}{9} = \dfrac{20}{36}$.
\end{exercise}

\begin{exercise}[$\star$]\label{exo:g7:fractions:2}
Sederhanakan sedapat mungkin:
$\dfrac{12}{16}$; \quad $\dfrac{30}{45}$; \quad $\dfrac{27}{36}$;
\quad $\dfrac{48}{12}$.
\end{exercise}

\begin{exercise}[$\star$]\label{exo:g7:fractions:3}
Apakah $\dfrac{15}{35}$ dan $\dfrac{21}{49}$ senilai? Dan
$\dfrac{16}{28}$ dan $\dfrac{20}{36}$? Berilah alasan dengan menyederhanakannya.
\end{exercise}

\begin{exercise}[$\star$]\label{exo:g7:fractions:4}
Bandingkanlah (tulislah penalarannya, bukan sekadar jawabannya):
\[
\frac{7}{11} \text{ dan } \frac{9}{11}; \qquad
\frac{5}{6} \text{ dan } \frac{11}{18}; \qquad
\frac{13}{12} \text{ dan } \frac{19}{20}.
\]
\end{exercise}

\begin{exercise}[$\star$]\label{exo:g7:fractions:5}
Hitunglah lalu sederhanakan kalau bisa:
\[
\frac{5}{9} + \frac{2}{9}, \qquad
\frac{11}{7} - \frac{4}{7}, \qquad
\frac{5}{8} + \frac{7}{8} .
\]
\end{exercise}

\begin{exercise}[$\star$]\label{exo:g7:fractions:6}
Hitunglah dengan \cref{met:g7:fractions:add}:
\[
\frac{2}{3} + \frac{5}{6}, \qquad
\frac{7}{10} - \frac{2}{5}, \qquad
\frac{3}{4} + \frac{5}{12} .
\]
\end{exercise}

\begin{exercise}[$\star$]\label{exo:g7:fractions:7}
Hitunglah:
\[
3 - \frac{5}{4}, \qquad
1 + \frac{3}{8}, \qquad
2 - \frac{7}{6} .
\]
\end{exercise}

\begin{exercise}[$\star$]\label{exo:g7:fractions:8}
Hitunglah:
\[
5 \times \frac{3}{10}, \qquad
\frac{7}{8} \times 4, \qquad
\frac{2}{3} \text{ dari } 45 .
\]
\end{exercise}

\begin{exercise}[$\star\star$]\label{exo:g7:fractions:9}
Adi memakan $\frac14$ bagian sebuah pizza dan Yuli $\frac38$ bagian
pizza yang sama. Berapa bagian pizza yang mereka makan bersama?
Berapa bagian yang tersisa?
\end{exercise}

\begin{exercise}[$\star\star$]\label{exo:g7:fractions:10}
Sebuah botol memuat $\frac34$ L jus. Lia menuangkannya $\frac16$ L
sebanyak dua kali. Berapa jus yang tersisa? (\omterm{def:g4:fractions:def}{Penyebut} bersamanya: $12$.)
\end{exercise}

\begin{exercise}[$\star\star$]\label{exo:g7:fractions:11}
Jelaskan, dengan alasan pada \cref{ex:g7:fractions:trap}, mengapa
$\frac35 + \frac12$ tidak mungkin sama dengan $\frac{4}{7}$, tanpa
menghitung \omterm{def:g1:addition:def}{jumlahnya} yang benar.
\end{exercise}

\begin{exercise}[$\star\star\star$]\label{exo:g7:fractions:12}
Hitunglah \omterm{def:g1:addition:def}{jumlah}
\[
\frac12 + \frac14 + \frac18 + \frac{1}{16}
\]
selangkah demi selangkah (dari kiri ke kanan). Amatilah pola \omterm{def:g1:addition:def}{jumlah}
sebagiannya: bilangan apa yang mereka dekati ketika suku
$\frac{1}{32}, \frac{1}{64}, \dots$ makin banyak ditambahkan?
\end{exercise}

\section{Soal: Pecahan yang bersembunyi dalam sebirama musik}

\begin{problem}[{Soal akhir pekan --- nilai nada adalah pecahan yang
harus berjumlah tepat, membilang irama menyembunyikan barisan yang
terkenal, dan keganjilan menguasai ketukan Balkan}]
\label{pb:g7:fractions:1}
Musik ditulis dengan \omterm{def:g6:fractions:def}{pecahan}. Sebuah \emph{nada penuh} berlangsung $1$
birama (dalam birama ``empat-empat'' yang biasa); \emph{nada setengah}
berlangsung $\frac12$; lalu datang \emph{nada \omterm{def:g3:division:half}{seperempat}} $\frac14$,
\emph{nada seperdelapan} $\frac18$ dan \emph{nada seperenam belas}
$\frac{1}{16}$. Satu hukum yang tidak boleh dilanggar: \emph{lama nada
di dalam satu birama harus berjumlah tepat sebesar biramanya} ---
yaitu $1$ pada birama empat-empat. Setiap pertanyaan di bawah adalah
hitungan bab ini (\cref{thm:g7:fractions:add},
\cref{met:g7:fractions:add}) yang digubah menjadi musik; di tengah
jalan kamu akan bertemu barisan pembilangan yang ditemukan para
cendekiawan India seribu tahun yang lalu.

\medskip
\noindent\textbf{Bagian I --- Nada yang harus berjumlah pas.}
\begin{enumerate}
  \item Periksalah bahwa setiap birama berikut sah pada birama
        empat-empat: (a) setengah $+$ \omterm{def:g3:division:half}{seperempat} $+$ \omterm{def:g3:division:half}{seperempat}; (b)
        \omterm{def:g3:division:half}{seperempat} $+$ seperdelapan $+$ seperdelapan $+$ setengah.
  \item Seorang penyalin menulis: setengah $+$ \omterm{def:g3:division:half}{seperempat} $+$
        seperdelapan. Berapa bagian birama yang tertulis, dan nada
        tunggal yang mana yang melengkapinya?
  \item Berapa nada seperenam belas yang mengisi satu birama penuh?
        Berapa seperenam belas yang menjadi satu nada seperdelapan?
  \item Sebuah \emph{titik} sesudah nada menambahkan setengah nilai
        nada itu: nada setengah bertitik bernilai $\frac12 + \frac14$,
        nada \omterm{def:g3:division:half}{seperempat} bertitik bernilai $\frac14 + \frac18$.
        Hitunglah kedua nilainya.
  \item Sebuah waltz ditulis pada birama tiga-empat: setiap biramanya
        harus berjumlah $\frac34$. Periksalah bahwa satu nada setengah
        bertitik mengisi satu birama waltz, lalu tulislah dua birama
        waltz sah yang lain yang hanya memakai nada \omterm{def:g3:division:half}{seperempat} dan
        seperdelapan.
\end{enumerate}

\medskip
\noindent\textbf{Bagian II --- Bengkel sang penyalin.}
\begin{enumerate}[resume]
  \item Sebuah birama empat-empat memuat: \omterm{def:g3:division:half}{seperempat} bertitik $+$
        seperdelapan $+$ \omterm{def:g3:division:half}{seperempat}. Apa yang kurang?
  \item Mana yang lebih lama, seperdelapan bertitik atau nada
        \omterm{def:g3:division:half}{seperempat}? Bandingkan lewat seperenam belas.
  \item Sebuah \emph{ikatan} merekatkan beberapa lama nada menjadi satu
        bunyi yang panjang: nada setengah yang diikat pada seperdelapan
        berlangsung berapa lama? Berapa seperenam belas itu?
  \item Nada sebuah birama berjumlah $\frac{11}{16}$; sisanya berupa
        kesenyapan (\emph{tanda istirahat}). Berapa lama tanda
        istirahat itu?
  \item Seorang penabuh mengisi satu nada \omterm{def:g3:division:half}{seperempat} ($\frac{4}{16}$
        biramanya) hanya dengan seperdelapan ($\frac{2}{16}$) dan
        seperenam belas ($\frac{1}{16}$), dengan urutan yang
        berpengaruh: misalnya
        seperdelapan--seperenam belas--seperenam belas, atau
        seperenam belas--seperdelapan--seperenam belas. Periksalah
        bahwa kedua contohnya sah, lalu sebutkan \emph{semua} irama
        yang mungkin. Ada berapa?
\end{enumerate}

\medskip
\noindent\textbf{Bagian III --- Barisan berumur seribu tahun, dan
sebuah birama \omterm{def:g3:numbers:evenodd}{ganjil}.}
\begin{enumerate}[resume]
  \item Mari kita bilang irama sang penabuh untuk nada yang lebih
        panjang. Setiap irama berakhir dengan seperenam belas atau
        dengan seperdelapan; jelaskan mengapa hal itu memberi aturan:
        (irama yang mengisi $n$ seperenam belas) $=$ (irama yang
        mengisi $n - 1$) $+$ (irama yang mengisi $n - 2$). Mulai dari
        $1$ irama untuk satu seperenam belas dan $2$ untuk dua,
        bangunlah tabelnya sampai $n = 8$: berapa irama yang mengisi
        satu nada setengah?
  \item Bilangan yang kamu temukan --- $1, 2, 3, 5, 8, 13, 21, 34$
        --- diterbitkan oleh cendekiawan India Hemachandra sekitar
        tahun 1150, ketika membilang irama semacam itu, dua generasi
        sebelum Fibonacci menuliskannya di Italia. Nyatakan pola yang
        mendefinisikannya dalam satu kalimat, lalu periksalah pada
        tiga entri terakhir tabelmu.
  \item Gaya ``swing'' mengganti dua seperdelapan yang sama dengan
        pasangan panjang--pendek: $\frac16$ lalu $\frac{1}{12}$.
        Buktikan bahwa penukaran itu sah (lama seluruhnya sama),
        dengan memakai \omterm{def:g4:fractions:def}{penyebut} bersama $12$.
  \item Sebuah \emph{triol} memampatkan tiga nada yang sama ke dalam
        satu nada \omterm{def:g3:division:half}{seperempat}. Berapa lama setiap notnya (periksalah
        jawabanmu dengan mengalikannya dengan $3$,
        \cref{prop:g7:fractions:mult})? Pertanyaan yang sama untuk
        tiga nada yang sama yang mengisi satu nada setengah.
  \item Musik tari Balkan memakai birama tujuh-delapan: setiap
        biramanya berjumlah $\frac78$. Buktikan bahwa \omterm{def:g3:division:half}{seperempat}
        bertitik $+$ \omterm{def:g3:division:half}{seperempat} $+$ \omterm{def:g3:division:half}{seperempat} mengisinya. Lalu
        jelaskan mengapa \emph{tidak ada} gabungan nada setengah dan
        nada \omterm{def:g3:division:half}{seperempat} saja, sebanyak apa pun, yang dapat mengisi
        birama tujuh-delapan. (Bilanglah dalam seperdelapan lalu
        pikirkan \omterm{def:g3:numbers:evenodd}{genap} dan \omterm{def:g3:numbers:evenodd}{ganjil}.)

\end{enumerate}
\end{problem}
